После распознавания текста средствами EasyOCR выполнялась дополнительная постобработка результата. Ее основная цель заключалась в уменьшении числа ошибок, возникающих на этапе OCR, и повышении читаемости итогового текста. В работе использовались два варианта постобработки: эвристическая нормализация и нормализация с применением языковой модели YandexGPT 5 Lite.

Эвристическая постобработка основана на наборе заранее заданных правил очистки и исправления текста. На этом этапе из результата OCR удалялись очевидные артефакты распознавания, нормализовались пробелы, устранялись лишние разрывы строк, приводились к единообразному виду знаки пунктуации и исправлялись типичные замены символов. К числу таких ошибок относятся подмена схожих по начертанию символов, появление лишних символов, пропуски разделителей между словами и другие характерные дефекты OCR.

Преимущество эвристического подхода заключается в его простоте, высокой скорости работы и полной предсказуемости результата. Однако его эффективность ограничена набором заранее предусмотренных правил. Если ошибка распознавания зависит от контекста или требует понимания смысла фразы, эвристические методы часто оказываются недостаточными.

Второй вариант постобработки основан на применении языковой модели YandexGPT 5 Lite. В этом случае распознанный OCR текст передавался в модель вместе с инструкцией исправить ошибки распознавания, сохранив исходный смысл, структуру и стиль фрагмента. Такой подход позволяет учитывать контекст соседних слов и восстанавливать более естественные текстовые конструкции даже в тех случаях, когда исходный результат OCR содержит несколько искажений одновременно.

Преимуществом использования языковой модели является более гибкая коррекция сложных ошибок, которые трудно устранить набором жестких правил. Особенно заметен эффект в случаях, когда OCR ошибается на уровне целых слов, нарушает границы слов или искажает отдельные буквенные последовательности. В таких ситуациях контекстная нормализация оказывается более эффективной, чем простая посимвольная замена.

При этом использование языковой модели требует аккуратной постановки инструкции. Важно ограничивать модель задачей исправления OCR-ошибок и не допускать произвольного переформулирования текста. Иначе существует риск появления изменений, которых не было в исходном документе. По этой причине постобработка с помощью LLM рассматривалась в работе именно как этап нормализации распознанного текста, а не как средство его свободного редактирования.